\documentclass[aps, amsmath,amssymb, reprint, twocolumn,notitlepage,superscriptaddress]{revtex4-2}

\setcitestyle{super}

\usepackage{graphicx}
\usepackage{dcolumn}
\usepackage{bm}
\usepackage{amsmath}
\usepackage{epstopdf}
\usepackage{epsfig}
\usepackage{subfigure}
\usepackage{fancyvrb}
\usepackage{mathbbol}
\usepackage{float}
\usepackage{tikz}

% REMOVED global "force to top" to allow figures to stay near text
% \makeatletter
% \setlength{\@fptop}{0pt}
% \makeatother

\newcommand{\br}{\mathbf{r}}
\newcommand{\bk}{\mathbf{k}}
\newcommand{\dd}[2]{\frac{\partial #1}{\partial #2}}

\begin{document}

\title{Is Pure Water a Mixture: Perspectives from Machine-Learning Classification Methods}
\author{Yiqi Yao}
\affiliation{Department of Chemistry, Harvey Mudd College, 301 Platt Boulevard, Claremont, California 91711}
\author{Diya Sanghi}
\affiliation{Department of Chemistry, Harvey Mudd College, 301 Platt Boulevard, Claremont, California 91711}
\author{Akanksha D. Chokshi}
\affiliation{Yale-NUS College, Singapore 138527, Singapore}
\author{Minwoo Choi}
\affiliation{Yale-NUS College, Singapore 138527, Singapore}
\author{Haiqin Wang}
\affiliation{Department of Chemistry, Harvey Mudd College, 301 Platt Boulevard, Claremont, California 91711}
\author{Bilin Zhuang}
\email{bzhuang@hmc.edu}
\affiliation{Department of Chemistry, Harvey Mudd College, 301 Platt Boulevard, Claremont, California 91711}
\affiliation{Yale-NUS College, Singapore 138527, Singapore}
\date{\today}

\begin{abstract}
Whether liquid water is a structural continuum or a dynamic mixture of two distinct local environments has remained a matter of serious debate for over a century. By applying unsupervised machine-learning clustering to molecular dynamics simulations of TIP4P/2005 and TIP5P water, we show that two structurally distinct local populations---locally favored tetrahedral structures (LFTS) and disordered normal-liquid structures (DNLS)---can be recovered without prior knowledge of the classification scheme. A systematic comparison of clustering algorithms operating on a five-dimensional order-parameter feature space ($q$, , LSI, $S_k$, $\zeta$) reveals that hybrid density-denoising pipelines yield the clearest two-state separation. Crucially, the cluster assignments are validated through an independent physical observable: the per-cluster oxygen--oxygen partial structure factor $S(k)$, computed via the Debye scattering equation directly from atomic coordinates. The LFTS cluster exhibits a first sharp diffraction peak at $k_{T1}\approx 0.81$, while the DNLS cluster peaks at $k_{D1} \approx 1.05$, in quantitative agreement with the theoretical predictions of the two-state model. Because the clustering and validation share no common descriptor, this agreement constitutes model-independent confirmation that the two-state signature in the structure factor reflects genuine structural heterogeneity in liquid water rather than an artifact of the classification scheme.
\end{abstract}

\maketitle

\section{Introduction}
The local structure of liquid water has been a subject of sustained scientific debate for over a century. Two fundamentally different pictures have competed to explain water's thermodynamic and dynamic anomalies---its density maximum at 4$^\circ$C, the rapid increase in isothermal compressibility upon supercooling, and the fragile-to-strong dynamic transition~\cite{gallo_water_2016}. Continuum models describe water through a broad unimodal distribution of local environments~\cite{narten_observed_1969, eisenberg_structure_2005}. Mixture models, dating back to R\"{o}ntgen's 1892 hypothesis of coexisting ``icebergs'' and a fluid ``sea''~\cite{rontgen-1892}, instead posit two or more distinct structural populations whose relative fractions shift with temperature and pressure. Despite a clear difference in physical picture, the lack of unambiguous experimental evidence has prevented convergence of this debate~\cite{narten_observed_1969, eisenberg_structure_2005, handle_supercooled_2017}.

Recent computational work by Shi and Tanaka has provided substantial support for a two-state description of liquid water~\cite{shi_direct_2020, russo_understanding_2014, shi_microscopic_2018}. In their framework, water is a dynamic mixture of locally favored tetrahedral structures (LFTS), stabilized by four hydrogen bonds with low density and high local symmetry, and disordered normal-liquid structures (DNLS), characterized by broken tetrahedral symmetry, higher density, and greater entropy. Using the structural descriptor $\zeta$---which quantifies translational order in the second coordination shell~\cite{russo_understanding_2014}---they demonstrated that the distribution $P(\zeta)$ is bimodal in several classical water models, and that the oxygen--oxygen partial structure factor $S(k)$ contains two overlapping peaks within the apparent first diffraction peak: a first sharp diffraction peak (FSDP) at $k_{T1} = kr_{\mathrm{OO}}/2\pi \approx 3/4$ arising from tetrahedral LFTS geometry, and an ordinary liquid peak at $k_{D1} \approx 1$ characteristic of DNLS~\cite{shi_direct_2020, shi-2019}. A methodological limitation persists in this prior work: the descriptor $\zeta$ that defines and classifies the two populations is also used to validate the classification through the $\zeta$-resolved structure factor $S(k, \zeta)$. This circularity leaves open the possibility that the observed two-state signatures are artifacts of the descriptor construction rather than intrinsic features of the liquid. An independent classification method that does not rely on $\zeta$-based thresholds would constitute a more rigorous test of the two-state hypothesis. In this work, unsupervised machine-learning clustering is applied to a five-dimensional order-parameter feature vector ($q$, $Q_6$, LSI, $S_k$, $\zeta$) computed from molecular dynamics trajectories of TIP4P/2005 and TIP5P water. Four clustering strategies are benchmarked---K-Means, Gaussian Mixture Models (GMM), and two hybrid pipelines combining density-based denoising (DBSCAN, HDBSCAN) with GMM---and the resulting cluster assignments are validated through the per-cluster oxygen--oxygen partial structure factor $S(k)$, computed via the Debye scattering equation directly from atomic coordinates. Because clustering is performed in order-parameter space and validation uses a reciprocal-space quantity derived independently from atomic coordinates, the two stages share no common descriptor, ensuring a non-circular test of cluster physical relevance.

\section{Method}
\subsection{Molecular Dynamics Simulations}
Classical molecular dynamics simulations are carried out for $N = 1024$ water molecules in the canonical (NVT) ensemble using the OpenMM simulation package~\cite{eastman_openmm_2015}. Two rigid, non-polarizable water models are employed: TIP4P/2005~\cite{abascal_general_2005} and TIP5P~\cite{mahoney_five-site_2000}, selected because Shi and Tanaka demonstrated that both produce strong tetrahedral ordering signatures with well-separated LFTS and DNLS populations in the structure factor~\cite{shi_direct_2020}. Simulations are conducted at seven temperatures from 0\,$^{\circ}$C to $-30$\,$^{\circ}$C in 5\,$^{\circ}$C increments, spanning conditions from a minority LFTS population to near-equal coexistence of both states. The Schottky temperatures---at which the two populations are approximately equal---are $T_{s=1/2} \approx 237.8$~K for TIP4P/2005 and $T_{s=1/2} \approx 255.5$~K for TIP5P~\cite{shi_direct_2020}. At each condition, 20 independent NVT production runs are performed following equilibration, yielding approximately $20{,}480$ molecular configurations per temperature per model. Trajectories are stored in .dcd and .pdb formats.

\subsection{Structural Order Parameters}
\label{sec:structural_order_parameter}
Five structural order parameters are computed for each molecule in every trajectory frame using the local oxygen--oxygen coordination environment. All parameters are calculated using MDTraj~\cite{mcgibbon_mdtraj_2015} for neighbor identification and custom implementations following the definitions below.
\begin{enumerate}
\item \textbf{Orientational Order Parameter, $q$}
Measures the extent to which a molecule and its four nearest neighbors form a tetrahedral arrangement~\cite{chau_new_1998, errington_relationship_2001}:
\begin{align}
q= 1 -\frac{3}{8}\sum_{j=1}^3 \sum_{k=j+1}^4{\left(\cos{\psi_{jk}}+ \frac{1}{3}\right)}^2
\end{align}
where $\psi_{jk}$ is the O--O--O angle between neighbors $j$ and $k$. Values range from $-3$ to $1$, with $q=1$ representing a perfect tetrahedron.
\item \textbf{Local Structure Index, LSI.}
The Local Structure Index quantifies the gap between the first and second hydration shells surrounding each molecule~\cite{shiratani_growth_1996}:
\begin{equation}
\text{LSI} = \frac{1}{n}\sum_{i=1}^{n}\left(\Delta_i - \bar{\Delta}\right)^{2}
\label{eq:LSI}
\end{equation}
where $\Delta_i = r_{i+1} - r_i$ is the distance gap between consecutively ordered oxygen neighbors ($r_1 < r_2 < \cdots < r_n$) within a cutoff of 3.7~\AA, and $\bar{\Delta}$ is the mean gap. A large LSI indicates a well-defined separation between shells, characteristic of tetrahedral ordering; a small LSI indicates a more uniform radial distribution of neighbors.
\item \textbf{Translational Tetrahedral Parameter, $S_k$}
Measures radial uniformity of the four nearest-neighbor distances~\cite{duboue-dijon_characterization_2015}:
\begin{align}
S_k= 1 -\frac{1}{3}\sum_{k=1}^4 \frac{(r_k - \Bar{r})^2}{4\Bar{r}^2}
\end{align}
where $r_k$ is the distance to the $k$th nearest neighbor and $\bar{r}$ is their mean. $S_k = 1$ for a perfect tetrahedron.
\item \textbf{Local Translational Order Parameter, $\zeta$}
Introduced by Russo and Tanaka~\cite{russo_understanding_2014}, $\zeta$ characterizes translational order in the second coordination shell:
\begin{equation}
\zeta = r_{\text{nearest non-HB}} - r_{\text{farthest HB}}
\label{eq:zeta}
\end{equation}
where $r_{\text{farthest HB}}$ is the distance to the most distant hydrogen-bonded neighbor and $r_{\text{nearest non-HB}}$ is the distance to the closest non-hydrogen-bonded neighbor. Large positive $\zeta$ indicates clear shell separation (LFTS); $\zeta \approx 0$ or negative indicates shell interpenetration (DNLS).
\end{enumerate}
All four parameters are computed per molecule per frame and stored as a feature matrix of dimensions $N_{\text{samples}} \times 5$, where $N_{\text{samples}} = N_{\text{frames}} \times N_{\text{molecules}} \times N_{\text{runs}}$.

\subsection{Clustering Methods}
\label{sec:clustering_algorithm}
All features are standardized to zero mean and unit variance prior to clustering. Six unsupervised methods are applied, falling into three categories~\cite{zheng-2018}:

\subsubsection{Baseline: K-Means and GMM}
K-Means~\cite{ikotun_k-means_2023} partitions the standardized feature space into $k=2$ disjoint clusters by iteratively assigning molecules to the nearest centroid. Centroids are initialized via K-Means++, with $n_{\text{init}} = 20$ restarts and a 300-iteration maximum. GMM~\cite{deprez-2024} provides a probabilistic alternative that is physically well-motivated: Tanaka's two-state framework predicts approximately Gaussian order-parameter distributions for each state~\cite{shi_direct_2020}. GMMs are fitted with $k=2$ full-covariance components using expectation-maximization~\cite{dempster_maximum_1977}; models with $k=2,3,4$ components are additionally evaluated using BIC and AIC to test for the presence of additional structural states.

\subsubsection{Hybrid Density--GMM Pipelines}
DBSCAN~\cite{ester_density-based_1996} identifies clusters as contiguous high-density regions and marks low-density boundary molecules as noise. In the hybrid pipelines, DBSCAN is applied solely as a denoising step: molecules in the low-density transition zone between the two structural states are identified and removed, and GMM is then fitted to the denoised subset. This two-stage approach combines the noise-identification capability of density-based algorithms with the physically motivated soft-clustering of GMM, yielding substantially higher silhouette score.

\subsection{Non-Circular Validation via Structure Factor}
\label{sec:structure_factor}
The central methodological contribution is a validation strategy that avoids the circularity of prior work, in which $\zeta$ is used both to define clusters and to validate them via $S(k, \zeta)$. Here, cluster labels from the ML pipeline are mapped back onto MD trajectories, and the oxygen--oxygen partial structure factor is computed independently for each cluster subpopulation using the Debye scattering equation~\cite{debye_zerstreuung_1915}:
\begin{equation}
S(k) = 1 + \frac{1}{N}\sum_{i=1}^{N}\sum_{\substack{j=1 \\ j\neq i}}^{N}\frac{\sin(k r_{ij})}{k r_{ij}}\,W(r_{ij})
\label{eq:debye}
\end{equation}
where $W(r_{ij}) = \sin(\pi r_{ij}/r_c)/(\pi r_{ij}/r_c)$ is a sinc window function that suppresses Gibbs ringing artifacts from the hard distance cutoff $r_c = 1.5$~nm. The wavenumber axis is normalized as $k r_{\text{OO}}/2\pi$, where $r_{\text{OO}} \approx 2.8$~\AA{} is the nearest-neighbor oxygen--oxygen distance, following the convention of Shi and Tanaka~\cite{shi_direct_2020}. If ML-assigned clusters correspond to physically distinct structural states, the LFTS cluster should exhibit an FSDP at $k_{T1} \approx 3/4$ and the DNLS cluster a peak at $k_{D1} \approx 1$---predictions entirely independent of the clustering step.

\section{Results and Discussion}
Algorithm benchmarking focuses on TIP4P/2005 water at $T = 253.15$~K ($-20\,^\circ$C), near the Schottky temperature $T_{s=1/2} \approx 237.8$~K where LFTS and DNLS populations coexist in appreciable proportions, providing a stringent test of each method's discriminative power.

\subsection{Clustering Algorithm Comparison}
\subsubsection{Baseline: K-Means and GMM}
K-Means ($k=2$) partitions the dataset of $N = 20{,}480$ molecular configurations into Cluster~0 (11{,}807 molecules, 57.7\%) and Cluster~1 (8{,}673 molecules, 42.3\%). The per-cluster order-parameter distributions and pairwise projections confirm that the two clusters align with LFTS and DNLS: $\zeta$ shows well-resolved bimodal peaks, and $q$ is enriched at high values in Cluster~0. The K-Means silhouette score is $0.2418$.

% --- FIGURE 1 k-mean baseline ---
\begin{figure}[ht!]
\centering
\includegraphics[width=0.95\textwidth]{Images/final_photos/fig1.png}
\subcaption{Per-cluster order-parameter distributions.}
\label{fig:kmeans_distributions}
\end{figure}

GMM yields qualitatively consistent populations (11{,}661 and 8{,}819 molecules) with a silhouette score of $0.2144$, modestly lower than K-Means due to the softer probabilistic boundaries. Together, the baseline methods confirm that a physically meaningful two-state partition is recoverable from the five-dimensional feature space without algorithmic specificity.

\subsubsection{Hybrid DBSCAN--GMM pipeline}
In the five-dimensional order-parameter space, LFTS and DNLS molecules occupy distinct dense regions separated by a sparsely populated transition zone. DBSCAN and HDBSCAN flag molecules in this low-density boundary as noise, providing a cleaner subset for GMM classification. 

% --- FIGURE 2 dbscan parameter ---
\begin{figure}[ht!]
\centering
\includegraphics[width=1.0\columnwidth]{Images/final_photos/fig2.png}
\caption{Parameter-selection heatmap for the DBSCAN denoising step in the hybrid DBSCAN--GMM pipeline. Stricter parameter choices (smaller $\varepsilon$ or larger MinPts) remove more boundary molecules and raise the silhouette score. Based on this analysis, $\varepsilon = 0.05$ and $\mathrm{MinPts} = 20$ are selected, removing $\sim$18.9\% of molecules as noise.}
\label{fig:dbscan_heatmap}
\end{figure}

Following the parameter sweep (Fig.~\ref{fig:dbscan_heatmap}), the DBSCAN--GMM pipeline removes approximately $22\%$ of molecules as noise and achieves a silhouette score of $0.2841$---a substantial improvement over both baselines. The sharpened cluster boundaries are evident in the per-cluster distributions and scatter plots (Fig.~\ref{fig:dbscan_gmm_combined}). The HDBSCAN--GMM pipeline removes $15.2\%$ as noise and achieves a comparable silhouette score of $0.2572$.

% --- FIGURE 3 MOVED HERE ---
\begin{figure}[ht!]
\centering
\begin{minipage}[t]{1.0\columnwidth}
\centering
\includegraphics[width=0.95\textwidth]{Images/final_photos/fig3.png}
\subcaption{Per-cluster order-parameter distributions.}
\label{fig:dbscan_gmm_distributions}
\end{minipage}
\caption{DBSCAN--GMM hybrid pipeline applied to TIP4P/2005 water at $T = 253.15$~K. Grey points indicate molecules removed as noise by DBSCAN; blue and orange denote LFTS and DNLS. The removal of the grey noise band visibly sharpens the boundary between structural populations.}
\label{fig:dbscan_gmm_combined}
\end{figure}

\subsection{Structure Factor Validation}
\label{sec:structure_factor_results}
The DBSCAN--GMM pipeline, which achieved the highest silhouette score, is used for structure factor validation. The per-cluster $S(k)$ curves (Fig.~\ref{fig:sk_per_cluster}) exhibit a clear shift in the position of the first diffraction peak: the LFTS cluster (Cluster~0) peaks at $k_{T1} \approx 0.81$ and the DNLS cluster (Cluster~1) peaks at $k_{D1} \approx 1.05$, in close quantitative agreement with the theoretical predictions $k_{T1} \approx 3/4$ and $k_{D1} \approx 1$ of Shi and Tanaka~\cite{shi_direct_2020}. 

% --- FIGURE 4 MOVED HERE ---
\begin{figure}[ht!]
\centering
\includegraphics[width=1.0\columnwidth]{Images/structure_factor_per_cluster_norm_tip4p2005_T-20.0.png}
\caption{Per-cluster oxygen--oxygen partial structure factors $S(k)$ for the DBSCAN--GMM pipeline. The LFTS cluster (blue) displays its primary peak at $k_{T1} \approx 0.81$, consistent with the FSDP of tetrahedral materials. The DNLS cluster (orange) peaks at $k_{D1} \approx 1.05$. This peak separation constitutes a non-circular validation of the two-state classification.}
\label{fig:sk_per_cluster}
\end{figure}

The total structure factor shows a single broad peak between $k_{T1}$ and $k_{D1}$, confirming that the apparent first diffraction peak of liquid water is a superposition of two overlapping contributions from structurally distinct subpopulations. The $\zeta$-resolved structure factor $S(k, \zeta)$ (Fig.~\ref{fig:sk_contour}) provides further confirmation: two distinct ridges are visible in separate $\zeta$ domains, with the FSDP at $k_{T1}$ concentrated in the high-$\zeta$ (LFTS) region and the ordinary peak at $k_{D1}$ dominating the low-$\zeta$ (DNLS) region, in quantitative agreement with the surfaces reported by Shi and Tanaka (cf. their Fig.~2D--E)~\cite{shi_direct_2020}. 

% --- FIGURE 5 MOVED HERE ---
\begin{figure}[ht!]
\centering
\includegraphics[width=1.0\columnwidth]{Images/2d_sk_zeta_combined_tip4p2005_T-20.0.png}
\caption{Two-dimensional contour map of the $\zeta$-resolved structure factor $S(k, \zeta)$. Two distinct intensity ridges are visible in separated $\zeta$ domains: high $\zeta$ (LFTS) corresponds to the FSDP, while low $\zeta$ (DNLS) corresponds to the ordinary liquid peak.}
\label{fig:sk_contour}
\end{figure}

The complementary three-dimensional $S(k, \zeta)$ surfaces per cluster (Fig.~\ref{fig:cluster_3d}) further illustrate the inverse peak structure---each cluster's $S(k)$ maximum corresponds to the other's minimum---providing non-circular confirmation of two physically distinct structural states.

% --- FIGURE 6 MOVED HERE ---
\begin{figure}[ht!]
\centering
\begin{minipage}[t]{0.48\textwidth}
\centering
\includegraphics[width=\textwidth]{Images/cluster0.png}
\end{minipage}%
\hfill
\begin{minipage}[t]{0.48\textwidth}
\centering
\includegraphics[width=\textwidth]{Images/cluster1.png}
\end{minipage}
\caption{Three-dimensional $S(k, \zeta)$ surfaces for the LFTS cluster (left) and DNLS cluster (right). Each cluster's $S(k)$ maximum corresponds to the other's minimum, resolving two physically distinct states.}
\label{fig:cluster_3d}
\end{figure}

\subsection{Robustness and Generality}

\subsubsection{Temperature Dependence}
The DBSCAN--GMM pipeline is applied to TIP4P/2005 water at five temperatures from $+10\,^\circ$C to $-30\,^\circ$C (Fig.~\ref{fig:sk_multitemp}). In the LFTS cluster, the FSDP at $k_{T1}$ intensifies monotonically upon cooling, consistent with the thermodynamic prediction that the LFTS fraction $s$ grows toward the low-density amorphous limit~\cite{shi-2019}. 

% --- FIGURE 7 MOVED HERE ---
\begin{figure}[ht!]
\centering
\includegraphics[width=1.0\columnwidth]{Images/fig_c3_sk_multitemp.png}
\caption{Temperature dependence of per-cluster $S(k)$ curves for TIP4P/2005 water. In the LFTS cluster (left), the first peak near $k_{T1} \approx 3/4$ intensifies monotonically upon cooling. The DNLS cluster (right) displays weaker temperature dependence.}
\label{fig:sk_multitemp}
\end{figure}

The DNLS cluster shows weaker temperature dependence, with its peak near $k_{D1}$ remaining stable across the full range. These trends mirror the behavior reported by Shi and Tanaka~\cite{shi_direct_2020} and confirm that the two-state classification is thermodynamically consistent across the supercooled regime.

\subsubsection{Optimal Cluster Number}
GMMs with $k = 1, 2, 3, 4$ components are fitted to TIP4P/2005 water at $-20\,^\circ$C and evaluated using BIC, AIC, and silhouette score (Fig.~\ref{fig:cluster_number}). BIC and AIC decrease monotonically, with the largest single drop between $k=1$ and $k=2$, confirming structural bimodality but without identifying a clear minimum---an expected result for thermally broadened overlapping populations. The silhouette score resolves this ambiguity: it peaks at $k=2$ ($0.163$) and declines thereafter, indicating that additional components fragment the physical populations rather than resolving genuinely distinct states.

% --- FIGURE 8 MOVED HERE ---
\begin{figure}[ht!]
\centering
\includegraphics[width=1.0\columnwidth]{Images/fig_c4_cluster_number.png}
\caption{Model selection analysis for the number of GMM components $k$. While BIC/AIC (top) decrease monotonically, the silhouette score (bottom) peaks at $k=2$, establishing it as the most parsimonious description.}
\label{fig:cluster_number}
\end{figure}

\section{Conclusion}
\label{sec:conclusion}
Unsupervised machine learning independently recovers two structurally distinct local populations in supercooled water, corresponding to the locally favored tetrahedral structures (LFTS) and disordered normal-liquid structures (DNLS) of Tanaka's two-state framework. Among the four clustering strategies benchmarked on TIP4P/2005 water at $T = 253.15$~K, the DBSCAN--GMM hybrid achieves the highest silhouette score ($0.2841$), outperforming K-Means ($0.2418$) and standalone GMM ($0.2144$) by effectively removing ambiguous transition-state molecules prior to probabilistic classification. The central result is the non-circular validation of the ML-assigned clusters via the oxygen--oxygen partial structure factor. The LFTS cluster exhibits an FSDP at $k_{T1} \approx 0.81$ and the DNLS cluster peaks at $k_{D1} \approx 1.05$, in quantitative agreement with the predictions of Shi and Tanaka~\cite{shi_direct_2020}. Because the clustering operates in order-parameter space and the validation uses a reciprocal-space observable derived independently from atomic coordinates, this agreement constitutes model-independent confirmation that the two-state signature in the structure factor is not an artifact of the classification scheme.

\newpage
\section*{Bibliography}
\bibliographystyle{apsrev4-2}
\bibliography{bib}
\end{document}
